\section{Conclusion}

The relaxed program has a proved conditional conclusion. For fixed
$1/3\lt\alpha\lt1/2$ and sufficiently large $Q$,

\begin{equation*}
a_Q(n)=1
\Longrightarrow
n\text{ is prime and }\Omega(n+2)\le2.
\end{equation*}

The divisor local factor property gives the exact one-divisor comparison. Wheel-sharing divisors
vanish; coprime divisors have

\begin{equation*}
\begin{aligned}
\mathcal N_m[L,U)
&=\rho(m)\ell_m+E_m[L,U),
&&[\text{Exact Interval Decomposition}]\\
|E_m[L,U)|&\le R-1,
&&[\text{Periodic Boundary}]\\
\rho_{Q,z}
&=
\frac12
\prod_{2\lt p\lt z}\left(1-\frac2p\right)
\prod_{z\le p\lt Q}\left(1-\frac1p\right).
&&[\text{Nested-Wheel Specialization}]
\end{aligned}
\end{equation*}

The cofactor progression discrepancy property identifies the natural pre-sieved Type-I remainder exactly:

\begin{equation*}
\begin{aligned}
r_d(I)
&=A_d(I)-\frac{A_1(I)}{\varphi(d)}
&&[\text{Shifted Divisor Discrepancy}]\\
&=\pi(I;d,-2)-\frac{\pi(I)}{\varphi(d)}.
&&[\text{Square-Safe Prime Identity}]
\end{aligned}
\end{equation*}

The bilinear character obstruction property gives the exact nonprincipal bilinear spectrum

\begin{equation*}
\mathbf1_{mn\equiv-2\ (\mathrm{mod}\ d)}
-\frac1{\varphi(d)}
=
\frac1{\varphi(d)}
\sum_{\chi\ne\chi_0}
\overline{\chi(-2)}\chi(m)\chi(n),
\end{equation*}

and the modulo-$3$ character refutes scalar-only centering:

\begin{equation*}
\left|
\sum_{m,n\in G_W}\chi_3(m)\chi_3(n)w(mn)
\right|
=
\sum_{m,n\in G_W}a(mn).
\end{equation*}

The next theorem must therefore add genuine averaged prime-distribution
information and use a locally adapted bilinear formulation. Even those two
estimates must still be inserted into a lower-bound almost-prime identity that
proves positivity. The article establishes the algebraic reductions and one
failed shortcut; it does not supply that final analytic argument, a new proof
of Chen's theorem, or a twin-prime result.
